% Prepared by the Part F agent; independently reviewed by root.
\subsection{The algebraic Hausdorff presentation and its comparison}

We specify the algebraic category, including its separation convention.
Its objects are smooth connected hyperbolic curves over $\C$.
Use the stacks $\mathcal M_{g,[n]}$ of unordered punctures;
labelled or coloured variants use the corresponding finite covers.
Generators are paths in their \'etale fundamental groupoids,
actual finite \'etale maps, and actual puncture fillings.
An actual isomorphism is identified with its canonical stack path.
Besides the path-groupoid laws, impose the following relations.
For every finite chain of specified covers and fillings, let
$\mathcal H_w$ be its diagram stack. A path in $\mathcal H_w$
from $w_0$ to $w_1$, inducing source path $\alpha$ and target
path $\beta$, imposes
\[
                       \beta w_0=w_1\alpha.
\]
Two chains of actual algebraic maps with the same composite
are also identified. All these relations may occur inside
longer words.

The \emph{Hausdorff \'etale presentation} means the universal
Hausdorff topological category with profinite automorphism groups
realizing these data, continuously on each \'etale path groupoid.
This is a universal property in terms of algebraic stacks and
their \'etale paths alone. Equivalently, relations are imposed
in a closed congruence. The following proof establishes its
existence, so existence is not an extra assumption in the definition.

\begin{claim}{E.8}{Proved}
The Hausdorff \'etale presentation exists. After choices of
surface markings, it is equivalent to
$\widehat{\mathcal C}_{\rm diag}$ of E.6. In particular,
\[
 \operatorname{Hom}(S,T)=
 \coprod_{uM_S\in\operatorname{Hom}_{\mathcal C_{\rm diag}}(S,T)/M_S}
       \overline{S_u}\backslash\widehat M_S,
 \qquad S_u=\{a\in M_S:ua=u\}.
\]
Composition is continuous, Hom spaces are compact, and
$\operatorname{Aut}(S)=\widehat M_S$. The completed pure-cover
category embeds faithfully, and the labelled puncture-forgetting
operations agree with those induced by the algebraic moduli maps.
This proves the comparison and completion required in E3 for
the isotopy relation of E.5.
\end{claim}
\begin{proof}
First the diagram stacks used in the definition exist and are
of finite type. A filling is described by a configuration of
extra disjoint points, with the retained and forgotten sets
distinguished. For a covering use the finitely many smooth
Hurwitz stacks of the given degree and boundary ramification
partitions. Non-Galois maps are included by
\citet[Theorem~6.20]{E_BR2007}; its opened proof identifies
the local deformation of a smooth cover with that of its
marked target. Iterated fibre products over the intermediate
moduli stacks give $\mathcal H_w$. Hyperbolicity makes the
automorphism groups finite. We retain these stack automorphisms:
a Hurwitz map to target moduli need not be representable when
deck automorphisms occur, as the footnote to
\citet[Proposition~6.14]{E_BR2007} explicitly records.

\emph{The Betti presentation.}
Temporarily use ordinary paths in the analytic stacks.
Contractibility of Teichm\"uller space and its quotient
description identify the fundamental groupoid of each moduli
stack with the groupoid of oriented surface markings up to
isotopy \citep[\S2.1]{F_Schneps2003}.
An analytic diagram path, after local topological trivializations
over an interval, gives a diagram isotopy with its endpoint
transports. Equal algebraic composites give equal underlying
maps. These observations define a functor from the Betti
presentation to $\mathcal C_{\rm diag}$.

Conversely, choose a complex structure on each surface object.
For a topological word ending in $T$, keep its chosen target
structure and propagate that structure backwards: a covering
pulls it back, and a filling restricts it to the punctured open.
All factors become holomorphic, extend over their compactifications,
and are algebraic. Join the resulting source structure to its
chosen structure by a path of marked complex structures.
Contractibility of Teichm\"uller space makes the path with
this prescribed identity marking unique up to the relevant
homotopy.

Here is the compatibility with the relations, including the
family issue. A topological diagram isotopy can be represented
at the level of its paths and endpoint transports by smooth
finite-configuration paths and finite-cover transports. For
fillings this is ordinary approximation of paths in the open
configuration manifold, relative to the endpoints, followed
by isotopy extension. Covering paths are uniquely transported
after their starting cover is fixed. Work backwards from the
target: these configuration paths and covering transports
give a path in the analytic diagram stack. This asserts a
continuous stack path; it does not assert that an arbitrary
nonsmooth original isotopy is a holomorphic family in its
real parameter. The induced source path has the prescribed
identity marking, so the naturality relation gives the
required equality of Betti arrows.

If two words have the same actual composite, their pulled-back
source complex structures are the same. Both factorizations
are simultaneously holomorphic and have equal algebraic
composites. The stipulated relation therefore identifies them.
For composition, transport the first diagram along the marked
complex-structure path used for the second; its naturality
square identifies this construction with that of the concatenated
word. These arguments prove independence of the choices and
give inverse functors, up to marking isomorphisms. The Betti
presentation is therefore equivalent to $\mathcal C_{\rm diag}$.
No two-sided-orbit assertion is used; E.7 is still a diagram
path relation.

\emph{All \'etale path relations.}
Let $\mathcal A$ be any Hausdorff topological category with
profinite automorphism groups, and let
$F:\mathcal C_{\rm diag}\to\mathcal A$ be a functor.
By the Betti comparison this is a realization of the Betti
presentation. Each induced homomorphism from a moduli
fundamental group to an automorphism group of $\mathcal A$
extends uniquely to the full profinite completion. Choosing
one ordinary path to each endpoint extends this assertion
to path groupoids; changing the choice changes the formula
by an ordinary loop and leaves the extension unchanged.

Noohi's Theorem~20.4 and its following corollary identify
these completions with the \'etale fundamental groups
\citep{E_Noohi2005}. The finite-type complex Deligne--Mumford
stacks in question have locally path connected, semilocally
simply connected analytic stacks. Thus ordinary path torsors
are dense in the \'etale path torsors, component by component.
For a path $\gamma$ in any one $\mathcal H_w$, the two sides
of its naturality equation are continuous functions of
$\gamma$ into a Hausdorff Hom space of $\mathcal A$.
They agree on ordinary paths, hence on their closure.
The equality therefore holds for every \'etale path, for
every diagram stack, and after every insertion into a
longer word. The equal-actual-composite relations do not
change. Conversely restricting a continuous \'etale
realization to ordinary paths gives the original functor.
The two constructions are inverse, also on compatible
natural transformations.

\emph{Existence and separation.}
By E.6, $\widehat{\mathcal C}_{\rm diag}$ represents functors
from $\mathcal C_{\rm diag}$ to every such $\mathcal A$.
The preceding paragraph identifies this universal property
with that of the algebraic Hausdorff \'etale presentation.
It proves both existence and the asserted comparison.

We make the stabilizer conclusion explicit. Every Hausdorff
realization kills $\overline{S_u}$, since its orbit map is
continuous and kills $S_u$. Conversely, if
$a\notin\overline{S_u}$, the category of E.6 is itself
one realization of the algebraic presentation, by the
preceding construction. Its distinct cosets
$\overline{S_u}$ and $\overline{S_u}a$ separate $u$
from $ua$. No universal closed relation can therefore
identify them. This proves the exact closure formula
without a subgroup-separability hypothesis or an
interchange of completion with an ordinary quotient.

The remaining structural statements follow from E.6.
For a pure covering, defect zero excludes all genuine
fillings at every intermediate relation by E.5. Its
discrete stabilizers are the covering-only ones, and so
are their closures. This proves the asserted completed
embedding. Labelled forgetting maps and their composites
are included in the presentation by the literal
algebraic forgetting maps and their equality relations.
\end{proof}

\begin{claim}{E.9}{Proved}
For the degree-two cover followed by one filling
\[
 f_b(z)=z^2+b,\qquad b\in Y:=\mathbb P^1\setminus\{0,1,\infty\},
\]
\[
 X_b=\mathbb P^1\setminus
 \{0,\infty,\sqrt{-b},-\sqrt{-b},\sqrt{1-b},-\sqrt{1-b}\},
\]
fix the three labelled target punctures. The diagram component
has stack $Y\times B\mu_2$. Its source braid subgroup is
the image of
\[
 \phi:\pi_1(Y)\times C_2\longrightarrow M_{0,[6]}.
\]
This is the right stabilizer of the arrow in the indicated
ramification component. Its completed stabilizer is exactly
$\overline{\operatorname{im}\phi}$, equal to the image
of $\pi_1^{\mathrm{et}}(Y\times B\mu_2)$ in
$\widehat M_{0,[6]}$.
\end{claim}
\begin{proof}
The displayed map gives a finite \'etale double cover of
$Y\setminus\{b\}$, followed by its one-puncture filling
into $Y$. On compactifications the two branch values are
$b$ and $\infty$. Every double cover with this labelled
target and these branch values is isomorphic to the
displayed one, uniquely up to its deck involution.
The family exists globally over $Y$, so the corresponding
$\mu_2$-gerbe is trivial: its stack is $Y\times B\mu_2$.
The source family gives $\phi$, with the $C_2$ factor
supplied by the global deck involution.

The source complexity is four, the target complexity is
one, and the degree is two, so the defect is two. The
single omitted ramification point over the interior
critical value $b$ already contributes all of that
defect. There are no further omitted interior points.
Thus any factorization of this composite into coverings
and fillings must first omit $b$ on the target so that
the degree-two factor is \'etale. Filling the source
ramification point before that factor would make it
ramified on the open. Apart from homeomorphisms and
trivial factors, the cover-then-fill pattern is forced.
Its diagram loops, with fixed labelled target, therefore
give exactly the right stabilizer $\operatorname{im}\phi$.

By Noohi's comparison, the \'etale group is the profinite
completion of $\pi_1(Y)\times C_2$. Its continuous image
in $\widehat M_{0,[6]}$ is compact and hence closed, and
the original image is dense in it. It is consequently
exactly $\overline{\operatorname{im}\phi}$.
No injectivity of $\phi$ or separability of its image
is required. This concrete braid calculation is
consistent with the general stabilizer proof of E.8.
\end{proof}

The Hausdorff clause in E.8 is deliberate. We do not
assert that the naive abstract quotient of finite words
in profinite paths has Hausdorff Hom spaces before this
separation. E3 is answered by the stated Hausdorff
presentation, with the precise closed-coset completion
and its algebraic comparison.
